%==============================================================================
% Section 2: Homogeneous baseline
%==============================================================================
\section{Homogeneous Baseline}
\label{sec:baseline}

This section organises the homogeneous baseline of the four Thurston
geometries $\Tthree$, $\Nilthree$, $\Sthree$, $\Solthree$ and fixes the
comparison axis.  The aim is to align the mode language, the baseline
invariants, and the minimal/activated distinction of each geometry in a common
convention before the bulk/boundary response dictionaries of the subsequent
sections are introduced.  The selection rules that appear later are read off
from the baseline comparison fixed here.

\subsection{Four-Geometry Setup}
\label{sec:baseline:setup}

All four geometries compared in this paper share the common Euclidean product
structure $M^3\times\Sone$, but the structure constants of the
three-dimensional spatial part, the Weyl curvature, the KK splitting pattern,
and the frame rigidity differ from geometry to geometry.  Even with the same
minisuperspace variables, the stability of the effective potential, the mass
splitting of the vector sector, and the activation conditions of the
Chern--Simons channel therefore vary in a geometry-dependent way.

We first fix the homogeneous sector as a comparison baseline, on top of which
the first-layer bulk response and the second-layer boundary response are
stacked.  This makes the words \emph{activated}, \emph{protected}, and
\emph{minimal} read in subsequent sections not as descriptive labels but as
clear comparative judgements relative to the homogeneous background.

\subsection{Basic Comparison Table}
\label{sec:baseline:table}

The current baseline is summarised in the following table.

\begin{table}[H]
\centering
\begin{tabular}{lcccc}
\toprule
Indicator & $\Tthree$ & $\Nilthree$ & $\Sthree$ & $\Solthree$ \\
\midrule
CS direction count & 0 & 1 & 3 & 0 \\
KK Higgsing & none & uniaxial & triaxial & biaxial \\
EC slice minimum & absent & present & absent & present \\
Maxwell baseline & universal & universal & universal & universal \\
spin-2 rigidity & non-rigid & non-rigid & non-rigid & rigid \\
$C^2_\LC$ & $0$ & $4/(3R^4)$ & $0$ & $16/(3R^4)$ \\
\bottomrule
\end{tabular}
\caption{Four-geometry homogeneous baseline.}
\label{tab:baseline}
\end{table}

This table conveys two facts simultaneously.  First, $\Nilthree$ and $\Sthree$
are non-trivial activated geometries with respect to both the CS direction
count and the Higgsing pattern.  Second, $\Solthree$ is on the minimal side
in the CS direction count but possesses a non-vanishing Weyl curvature and
hence an EC slice-minimum branch isomorphic to that of $\Nilthree$.
Accordingly $\Solthree$ does not belong to the inert class of $\Tthree$; it is
a CS-inert but spin-0 active rigid geometry.  In the boundary layer this bulk
distinction is reinforced by the presence or absence of spectral branches
depending on the compact quotient and the spin structure.

\subsection{Maxwell Universality and the CS Direction-Count Law}
\label{sec:baseline:maxwell}

The first universal statement in the bulk layer is the geometry-independence
of the Maxwell coefficient.  Throughout this paper we adopt
\begin{equation}
  \KMxw = -\frac{L^2}{2r_0^4}
\end{equation}
as the baseline coefficient common to the four geometries.  This universality
keeps the skeleton of the main dictionary intact under inhomogeneous opening
and under either reading of the Euclidean circle.  Geometry dependence
therefore appears not in the Maxwell sector itself but in the CS-active
direction count, the Higgsing pattern, the EC slice-minimum structure, and
the defect response.

By contrast, Chern--Simons activation is geometry-dependent and the direction
count is set by the off-diagonal structure.  Activation occurs along one
direction in $\Nilthree$ and along three directions in $\Sthree$, while in
$\Solthree$ a self-referential cancellation forces it to zero.  The
distinction is decided not by ``whether structure constants exist'' but by
``whether the correction to $\tilde{F}_{ij}$ has its leg outside the index
pair $\{i,j\}$''.

The cancellation in $\Solthree$ arises because the underlying structure
constants form the symmetric pair
\begin{equation}
  C^1{}_{01} = +\frac{1}{R}, \qquad
  C^2{}_{02} = -\frac{1}{R},
\end{equation}
so that only same-leg self-coupling corrections remain.  Specifically the
correction $A_1\to\tilde{F}_{01}$ has its leg $1$ inside $\{0,1\}$ and is thus
self-referential, whereby the contributions $+1/R$ and $-1/R$ cancel
algebraically (off-diagonal CS rule, \S\ref{sec:rules:bulk} R2; details in
Appendix~\ref{app:E:E3}).  In this case the CS channel does not generate any
new parity-odd direction; only the biaxial Higgsing remains on top of the
Maxwell baseline.

\subsection{Mode Dictionary and Higgsing Pattern}
\label{sec:baseline:dict}

Read in the mode language, the differences among the four geometries in the
homogeneous sector concentrate in three sectors.  In the spin-0 sector, the
EC slice minimum on the $\eta=V=0$ section appears in $\Nilthree$ and
$\Solthree$.  Both share the same stationary radius and the same full
homogeneous Hessian criterion, but the slice potential and the radial Hessian
of $\Solthree$ are four times those of $\Nilthree$.  $\Tthree$ has a flat
zero slice and no isolated branch, while round $\Sthree$ has a non-zero
slope and no stationary point on the slice (verified in
Appendix~\ref{app:E:E5b}).  In the spin-1 sector the Higgsing pattern is
classified into none/uniaxial/triaxial/biaxial; this constitutes the principal
comparison axis of the bulk dictionary.  In the spin-2 sector only
$\Solthree$ exhibits frame rigidity; the other three geometries belong to the
non-rigid side.

In particular, the biaxial Higgsing of $\Solthree$ takes the form
\begin{equation}
  m^2(A_1) = m^2(A_2) = -\frac{L^2}{2r_0^4}, \qquad m^2(A_0) = 0.
\end{equation}
This follows from the modified field strength
$\tilde{F}_{01} = F_{01} + A_1/R$, which induces a mass term for $A_1$ and
$A_2$ with the same coefficient as $\KMxw$ (same-leg self-coupling), while
$A_0$ does not appear in $\tilde{F}$ and remains massless (KK derivation in
Appendix~\ref{app:E:E2}).  The fact that the two massive directions share the
same coefficient as the Maxwell baseline shows that $\Solthree$ is a distinct
geometry: it does not introduce a CS-active direction but carries an EC slice
branch and a rigid scaffold.

The four Higgsing types ``none/uniaxial/triaxial/biaxial'' are realised
exactly once each by $\Tthree, \Nilthree, \Sthree, \Solthree$, respectively,
as can be verified by direct computation of the KK mass matrices
(Appendix~\ref{app:E:E5}).

The role of this section is therefore to fix the baseline that supplies the
comparison axis of the response dictionary.  Most technical details---the
$\Solthree$ structure constants, the CS cancellation, and the frame
rigidity---are deferred to Appendices~\ref{app:A} and~\ref{app:E}.
